% Auto-generated by experiments/analysis/build_priorwork_table.py.
% Re-run that script to regenerate; do not edit by hand.
\begin{table}[t]
\centering
\caption{Cited prior-work numbers on EEG-based identification. Each row uses its source paper's protocol; no baseline is re-implemented under the AEP-Hybrid protocol of this paper. V4 results are in Table~\ref{tab:cross}.}
\label{tab:priorwork}
\begin{threeparttable}
\scriptsize
\setlength{\tabcolsep}{3pt}
\begin{tabular}{l l c c c}
\toprule
\textbf{Method} & \textbf{Dataset} & \textbf{Subj.} & \textbf{Metric} & \textbf{Score} \\
\midrule
MindID~\cite{zhang2017mindid}\tnote{1} & EEG-ID & 8 & Acc. & 0.984 \\
EEG-Triplet~\cite{yang2022eeg}\tnote{2} & PhysioNet MI & 109 & Acc. & 0.961 \\
BrainNet~\cite{fallahi2023brainnet}\tnote{3} & ERP-CORE & 40 & EER & 0.027 \\
AISEI v0~\cite{vu2026aisei}\tnote{4} & PhysioNet AEP & 20 & P@1 & 0.860 \\
\bottomrule
\end{tabular}
\begin{tablenotes}\scriptsize
\item[1] Within-session resting-state EEG; attention-RNN; no synthetic noise injection.
\item[2] Motor-imagery task; triplet loss on raw signals; session-mixed split.
\item[3] Siamese network on ERP epochs; cross-session evaluation; open-set verification on enrolled subjects only.
\item[4] Same dataset and same AEP-Hybrid protocol as this work; WaveNet-only denoiser; power-line noise.
\end{tablenotes}
\end{threeparttable}
\end{table}
